
% =====================================================================
% APÊNDICE — DECLARAÇÕES DE USO DE INTELIGÊNCIA ARTIFICIAL GENERATIVA (IAG)
% Conforme Portaria CNPq nº 2.664, de 6 de março de 2026
% (Institui a Política de Integridade na Atividade Científica do CNPq)
%
% Este arquivo foi pensado para ser incluído via \input{} ou \include{}
% em um relatório maior, dentro do ambiente de apêndices.
% Requer os pacotes: geometry (opcional), enumitem, array, xcolor (opcional)
% =====================================================================

\chapter{Declarações de Uso de Inteligência Artificial Generativa (IAG)}
\label{apendice_declaracao_iag}

\noindent
Em conformidade com a Portaria CNPq nº 2.664, de 6 de março de 2026, que institui a
Política de Integridade na Atividade Científica do CNPq e, em seu Art.~9º, inciso I,
alíneas ``c'' e ``d'', determina a obrigatoriedade de declaração do uso de ferramentas
de Inteligência Artificial Generativa (IAG) em qualquer fase do desenvolvimento da
pesquisa --- concepção, redação, análise de dados ou submissão ---, especificando a
ferramenta utilizada e a finalidade, apresentam-se a seguir as declarações relativas
ao uso de IAG neste trabalho.

\bigskip
\noindent
Ressalta-se que, nos termos da referida Portaria, é vedada a submissão de conteúdo
gerado por IAG como se fosse de autoria humana, sendo os autores integralmente
responsáveis pelo conteúdo final, inclusive por eventuais plágios ou imprecisões
geradas pela ferramenta (Art.~9º, I, ``d'' e ``f'').

% =====================================================================
% MODELO 1 — DECLARAÇÃO DE USO DE IAG PARA REVISÃO DE TEXTO
% =====================================================================

\section{Modelo de Declaração --- Revisão de Texto}
\label{ap:declaracao_revisao_texto}

\begin{center}
\textbf{\large DECLARAÇÃO DE USO DE INTELIGÊNCIA ARTIFICIAL GENERATIVA}\\
\textbf{(Revisão Textual e Linguística)}
\end{center}

\bigskip
\noindent
Eu, \underline{\hspace{7cm}} (nome completo do(a) autor(a) ou responsável), aluno da disciplina Práticas de Extensão I, II, III ou IV, DECLARO, para os devidos fins e em atendimento ao disposto na Portaria CNPq nº 2.664/2026, que fiz uso de ferramenta(s) de Inteligência Artificial Generativa (IAG) no desenvolvimento deste trabalho, conforme especificado abaixo:

\bigskip
\begin{tabular}{|p{4.5cm}|p{9.5cm}|}
\hline
\textbf{Ferramenta(s) utilizada(s)} & \underline{\hspace{9cm}} \\
& (ex.: ChatGPT, Claude, Grammarly, DeepL Write, etc. --- indicar nome e versão/modelo) \\
\hline
\textbf{Fase(s) da pesquisa em que houve uso} &
\begin{itemize}[leftmargin=*, nosep]
  \item[$\square$] Concepção
  \item[$\square$] Redação
  \item[$\square$] Análise de dados
  \item[$\square$] Submissão
\end{itemize} \\
\hline
\textbf{Finalidade do uso} & Revisão gramatical, ortográfica e de estilo do texto;
melhoria de clareza e coesão textual; correção de fluência da linguagem.
\emph{(Especificar detalhes adicionais, se houver)}: \underline{\hspace{9cm}} \\
\hline
\textbf{Trechos/seções do trabalho em que a ferramenta foi aplicada} &
\underline{\hspace{9cm}} \\
\hline
\end{tabular}

\bigskip
\noindent
Declaro, ainda, que:
\begin{enumerate}[label=\alph*)]
  \item o conteúdo intelectual, as análises, as interpretações e as conclusões
  apresentadas neste trabalho são de minha inteira autoria e responsabilidade;
  \item a ferramenta de IAG foi empregada exclusivamente como apoio à revisão
  linguística e estilística do texto, não tendo sido utilizada para geração de
  conteúdo científico, dados, resultados ou conclusões;
  \item todo o texto revisado por IAG foi integralmente analisado, verificado e
  validado por mim antes da submissão, assumindo integral responsabilidade por
  eventuais imprecisões, distorções de sentido ou plágios porventura introduzidos
  pela ferramenta;
  \item não foi submetido conteúdo gerado por IAG como se fosse de autoria humana,
  em conformidade com o Art.~9º, I, ``d'', da Portaria CNPq nº 2.664/2026.
\end{enumerate}

\bigskip
\noindent
Pinhais e data: \underline{\hspace{4cm}}, \underline{\hspace{1.5cm}} de
\underline{\hspace{3cm}} de \underline{\hspace{1.5cm}}.

\bigskip\bigskip
\noindent
\underline{\hspace{7cm}}\\
Assinatura do(a) autor(a) / responsável

% =====================================================================
% MODELO 2 — DECLARAÇÃO DE USO DE IAG PARA REVISÃO DE BIBLIOGRAFIA
% =====================================================================

\newpage
\section{Modelo de Declaração --- Revisão e Levantamento Bibliográfico}
\label{ap:declaracao_revisao_bibliografia}

\begin{center}
\textbf{\large DECLARAÇÃO DE USO DE INTELIGÊNCIA ARTIFICIAL GENERATIVA}\\
\textbf{(Levantamento, Revisão e Organização Bibliográfica)}
\end{center}

\bigskip
\noindent
Eu, \underline{\hspace{7cm}} (nome completo do(a) autor(a) ou responsável),
aluno da disciplina Práticas de Extensão I, II, III ou IV, DECLARO, para os devidos
fins e em atendimento ao disposto na Portaria CNPq nº 2.664/2026, que fiz uso de
ferramenta(s) de Inteligência Artificial Generativa (IAG) para apoio ao levantamento,
à revisão e/ou à organização de referências bibliográficas neste trabalho, conforme
especificado abaixo:

\bigskip
\begin{tabular}{|p{4.5cm}|p{9.5cm}|}
\hline
\textbf{Ferramenta(s) utilizada(s)} & \underline{\hspace{9cm}} \\
& (ex.: ChatGPT, Claude, Perplexity, Elicit, Consensus, Scite, etc. --- indicar nome
e versão/modelo) \\
\hline
\textbf{Fase(s) da pesquisa em que houve uso} &
\begin{itemize}[leftmargin=*, nosep]
  \item[$\square$] Concepção
  \item[$\square$] Redação
  \item[$\square$] Análise de dados
  \item[$\square$] Submissão
\end{itemize} \\
\hline
\textbf{Finalidade do uso} &
\begin{itemize}[leftmargin=*, nosep]
  \item[$\square$] Busca e sugestão de referências relacionadas ao tema
  \item[$\square$] Organização/formatação de referências (normas ABNT, APA, Vancouver etc.)
  \item[$\square$] Síntese ou resumo de artigos previamente selecionados pelo(a) autor(a)
  \item[$\square$] Verificação de completude/consistência de citações no texto
  \item[$\square$] Outra finalidade: \underline{\hspace{5cm}}
\end{itemize} \\
\hline
\end{tabular}

\bigskip
\noindent
Declaro, ainda, que:
\begin{enumerate}[label=\alph*)]
  \item todas as referências sugeridas ou organizadas com apoio da ferramenta de IAG
  foram individualmente verificadas quanto à sua existência, exatidão bibliográfica
  e pertinência ao tema, mediante consulta às fontes originais;
  \item nenhuma referência foi incluída no trabalho sem confirmação de sua
  autenticidade, sendo de minha exclusiva responsabilidade eventuais erros,
  omissões ou citações fabricadas (\emph{``alucinações''}) porventura geradas pela
  ferramenta;
  \item a seleção crítica, a interpretação e o uso das fontes bibliográficas no
  corpo do trabalho são de minha inteira autoria e responsabilidade;
  \item não foram inseridos, nas ferramentas de IAG utilizadas, projetos de pesquisa
  de terceiros ou dados de caráter sigiloso para fins de elaboração de pareceres,
  buscas ou análises, em conformidade com o Art.~9º, I, ``e'', da Portaria CNPq
  nº 2.664/2026.
\end{enumerate}

\bigskip
\noindent
Pinhais e data: \underline{\hspace{4cm}}, \underline{\hspace{1.5cm}} de
\underline{\hspace{3cm}} de \underline{\hspace{1.5cm}}.

\bigskip\bigskip
\noindent
\underline{\hspace{7cm}}\\
Assinatura do(a) autor(a) / responsável

% =====================================================================
% Nota de rodapé de embasamento legal (opcional — pode ser removida)
% =====================================================================
\bigskip
\noindent\rule{\textwidth}{0.4pt}
\footnotesize
\noindent
\textbf{Base legal:} Portaria CNPq nº 2.664, de 6 de março de 2026, Art.~9º, inciso I,
alíneas ``c'' (dever de declarar o uso de IAG, especificando ferramenta e finalidade),
``d'' (vedação à submissão de conteúdo de IAG como autoria humana), ``e'' (vedação à
inserção de projetos de pesquisa de terceiros em ferramentas de IAG) e ``f''
(responsabilidade integral do autor pelo conteúdo final, inclusive por plágios ou
imprecisões geradas pela IAG).
\normalsize
